\section{Experimental setup}
\label{sec:setup}

\paragraph{Models.} We evaluate \NModels{} models reachable through one gateway, spanning four providers and
several capability tiers: \texttt{gpt-6-astra}, \texttt{gpt-6-sol}, \texttt{gpt-5.6-sol}, \texttt{gpt-5.4-mini}
and \texttt{gpt-4.1} (OpenAI), \texttt{claude-opus-5.5} (Anthropic), \texttt{gemini-3.8-flash} (Google) and
\texttt{grok-4.7} (xAI). All models are used with provider-default
sampling and reasoning settings. Model identifiers are those exposed by the gateway at the time of the study
(September 2026). One preregistered model was withdrawn before publication; it is excluded from all analyses and
from the released data.

\paragraph{Minimal scaffold.} The scaffold is a standard function-calling loop \citep{yao2023react}: a fixed
system prompt (Appendix~\ref{app:prompts}), the task as the first user message, the tools of the task's
services plus the three general tools, and up to 40 tool calls. Parallel tool calls are executed in order.
If a turn contains no tool call, the scaffold asks the model once more to continue or finish; after three
such turns the episode ends. The same conversation is rendered to the chat-completions or the responses
protocol depending on the model, with reasoning items replayed where the protocol supports it.

\paragraph{Production harnesses.} We run three agent command-line interfaces non-interactively in an empty
temporary directory with \bench{} as their only MCP server: GitHub Copilot CLI 1.0.86, Hermes Agent 0.20.6 and
OpenAI Codex CLI 0.139. Each harness keeps its own system prompt, tool-call protocol, context management and
retry logic; we prepend the same operational preamble to the task (Appendix~\ref{app:prompts}), restrict the
harness to the sandbox's tools where it allows this, and give each episode a fresh harness home directory.
All harnesses reach their models through the same gateway as the scaffold, which lets us cross harnesses with
models.

\paragraph{Contract variants.} In the \emph{native} contract (Table~\ref{tab:services}) only billing and one
social platform accept idempotency keys. The counterfactual \emph{keys-everywhere} contract extends
Stripe-style key semantics to every non-idempotent write, including resumable batches, late commits and
redelivered requests, and documents the new parameter in each tool description. Nothing else changes.

\paragraph{Recovery conditions.} Two conditions change only the prompt: \emph{aware} adds five sentences of
reliability guidance to the system prompt, and \emph{reflect} inserts a generic reflection request after any
turn with a tool error \citep{shinn2023reflexion}. The remaining conditions wrap tool execution, as a framework or
middleware would, and never see ground truth: \emph{sdk-retry} transparently retries timeouts, 5xx and 429
responses up to three times with exponential backoff, mimicking common client defaults; \emph{rules} retries
only 429 and 503 responses; \emph{vbr} (verify-before-retry, after \citet{mansoor2026verified}) intercepts a
write whose identical predecessor returned an ambiguous error and runs its declared read-back first, sending
the write only if the effect is absent; \emph{wait~$\Delta$} (E5) additionally assumes an in-flight bound of
$\Delta$ seconds and defers that read-back until $\Delta$ seconds, plus the documented visibility lag, after the
original request was sent; and \emph{guard} adds four contract-driven components to vbr---automatic
idempotency keys where the contract accepts them, re-verification after the documented visibility lag,
blocking an unverifiable non-idempotent repeat with a request to escalate, and an annotation telling the
model that an outcome is unknown. Two reference conditions may read ground truth. The \emph{state oracle} is the
guard verifying against the current ground-truth state; it cannot see a request that is still in flight, so it
is \emph{not} an upper bound under late commits. The \emph{outcome oracle} also sees in-flight requests and
suppresses any re-issue of a write that will commit; it is the upper bound for any client-side policy, since only
redelivery on writes without key support remains beyond its reach. In harness runs the wrappers execute inside
the MCP server, so they apply to any harness without modification.

\paragraph{Experiments.} E1 crosses the \NModels{} models with the eleven fault modes of the original design (all but
the heavy-tailed late commit) and every focal write of two instances per template in the minimal scaffold, plus a
fault-free episode per instance (\NEOne{} episodes with the supplement below), and adds instances for the two
sparsest contract classes (partial batches and the unverifiable endpoint); \texttt{gpt-4.1} is served from a
separate low-rate quota and was run in its own low-concurrency lane with the identical design. E2 crosses three
models (\texttt{claude-opus-5.5}, \texttt{gpt-6-sol}, \texttt{gemini-3.8-flash}) with
nine conditions (vanilla, aware, reflect, sdk-retry, rules, vbr, guard and the two oracles) and the eight fault
modes most relevant to recovery, and repeats vanilla and guard under the keys-everywhere contract (E2k; \NETwo{}
episodes in total). E3 runs Copilot CLI and Hermes with \texttt{claude-opus-5.5}, \texttt{gpt-5.6-sol} and
\texttt{gemini-3.8-flash}, and Codex CLI (which speaks only the responses protocol) with \texttt{gpt-5.6-sol} and
\texttt{gpt-6-sol}, natively and with the guard in the MCP server, on the core fault modes and two focal writes per
template; the minimal scaffold is run on the identical design as the reference harness. E3k repeats the harness
design under the keys-everywhere contract with one model shared by all four harnesses (\texttt{gpt-5.6-sol};
\NEThree{} episodes for E3 and E3k). E4 probes robustness on two focal writes per template, compared with matched
E2 episodes: tool descriptions without consistency statements or with explicit non-idempotency warnings, three
paraphrases of the system prompt, an operator who never answers, and guard ablations (\NEFour{} episodes). E5 runs
the wait-then-verify family ($\Delta \in \{0, 60, 120, 300\}$\,s under the fixed delay; $\{0, 60, 300, 900, 3600\}$\,s
under the heavy-tailed delay) against vanilla, the guard, the outcome oracle and the keys-everywhere contract on the
three E2 models and all focal writes of instance 0 (\NEFive{} episodes). E6 removes the closing ``exactly once''
sentence from every instruction for five models on instance 0 and pairs each episode with its E1 counterpart
(\NESix{} episodes).

\paragraph{Statistics.} The world seed depends only on (template, instance, focal write, fault mode), so all
models, harnesses, conditions and instruction variants face identical worlds, and comparisons are paired.
Episodes that share a task template are correlated, and there are only twelve templates; a cluster bootstrap over
so few clusters is anti-conservative and degenerate for cells with no events. Per-cell 95\% intervals are therefore
Wilson intervals on a Kish effective sample size $n/(1+(\bar m-1)\rho)$, where $\bar m$ is the mean number of
episodes per template in the cell and $\rho$ the within-cell intra-class correlation by template, estimated once
per experiment. Regression analyses use a Bayesian binomial mixed model with random intercepts for template and
template$\times$instance (as preregistered for H1, fitted by variational Bayes) and, as a frequentist check,
generalized estimating equations clustered by template with bias-reduced (Mancl--DeRouen) standard errors and a
$t$ reference with $G-1$ degrees of freedom. Paired binary outcomes use exact McNemar tests with Holm correction
within each family; H4 uses the preregistered two one-sided tests (TOST) with a $\pm$5\,pp margin; observation
equivalence is assessed by comparing first-action agreement across matched worlds with agreement across independent
runs of the same world, and by a within-pair permutation test of the total variation distance. The Shapley
decomposition of McFadden's pseudo-$R^2$ is reported pooled, as preregistered, and separately for faults that an
immediate read-back resolves and faults it cannot. Episodes whose focal write was never attempted are excluded
from fault-conditional analyses and reported as trigger rates; episodes that end in a gateway error after eight
attempts are excluded and re-run. Hypotheses, metrics and analyses were preregistered before E1; E2k, E3k, E5, E6,
the outcome oracle, the stratified decomposition and the revised interval method were added afterwards and are
exploratory, and every deviation from the preregistration is listed in the released materials.
